\documentclass[]{../../../../tex/classes/styledArticle}
\usepackage{../../../../tex/packages/hebrewSupport}

\author{שחר פרץ}
\title{\textit{היסטוריה 46}}
\begin{document}
	\maketitle
	מלחמת סיני הסתיימה בהסכמים מסויימים, לגבי פירוז שטחים בסיני, ופתיחת מצרי טירן. ברגע שמצרי טירן נסגרו נחסם נמל אילת. מדינת ישראל יזמה את מלחמת ששת הימים כתוצאה מהאירועים האלו ועוד אירועים שפורטו בשיעור הקודם, בה כל חצי האי סיני, עיו''ש, עזה והגולן נכבש. מדינת ישראל שילשה את שטחה. בתקופת המלחמה נכנסו לישראל כמיליון פלסטינים שרובם חיים בעיו''ש. 
	
	החל הויכוח ברמה הלאומית על החזרת שטחים. התחושה של החרדה הקיומית התחלפה באופוריה, זלזול  באויב ושאננות. חל גידול ניכר בעולים לארץ. מצבה הכלכלי מדינת ישראל השתפר. מערכת היחסים עם ברה''מ וצרפת עם ישראל התנתקה (צרפת היא זו שסייע למדינת ישראל לבנות את מפעל הטקסטיל למטרות שלום בדימונה). האו''ם שמום קובע את החלטה 242 בנוגע לנסיגה משטחים תמורת שלום. 
	
	ההשפעות של מלחמת ששת הימים קשורים באופן ישיר לגורמים למלחמת יום הכיפורים. כאשר מדברים על נסיבתיות, נדבר על שלושה דברים בהקשר של מלחמת יום הכיפורים: 
	\begin{multicols}{2}
		\begin{itemize}
			\item גורמים/סיבות למלחמה (קשור \textit{רק} למדינות ערב)
			\item גורמים להפתעה
			\item נסיבות
			\item תוצאות
			\item השפעות
		\end{itemize}
	\end{multicols}
	
	בנוגע לגורמים למלחמה: 
	\begin{itemize}
		\item רצון להחזיר את הכבוד, ולשפר את מעמד המנהיגים בתוך המדינה
		\item סאדאת רוצה פתרון מדיני, שטחים תמורת שלום. הרצון הזה לא מתממש. סאדאת היה מנהיג יחסית חלש לעומת נאצר שקדם אותו. 
		\item הרצון של מדינות ערב להחזיר את השטחים שנכבשו בששת הימים
	\end{itemize}
	
	בהקשר של מצריים, מ־1969 החלו דיבורים על שלום על ישראל. כל עוד נאצר היה במצרים ההצעות נדחו משתי הצדדים, אך כאשר סאדאת עלה המצב התחלף, בגלל כל מני עניינים פנימיים. גולדה מאיר סירבה למסע ומתן האמריקקי. 
	
	רק אחרי 50 שנה התחילו להפתח פרוטוקולים בנוגע למלחמה. כתוב במגוון ספרים ומקורות שגולדה לא הסכימה לשלום. מהפרוטוקולים שנפתחו בשנים האחרונות, הסיבה העיקרית לכך הייתה הדרישה המצרית שהמסע ומתן יהיה עם תנאים מוקדמים, כלומר עוד מתחילת המסע ומתן ישראל תצא משטחים. גולדה הסכימה להכנס למסע ומתן, ללא תנאים מקדימים. 
	
	סיני היא איזור משמעותי למצריים – יש בה מאגרי נפט, הוא שטח נרחב, וכו'. 
	
	מדינות ערב הפיקו לקחים ממלחמת ששת הימים. 
	\begin{itemize}
		\item לחתור ממלחמה במספר חזיתות כדי למנוע ריכוז כוחות
		\item השענות על אחת המעצמות (כנראה ברה''מ)
		\item תגבור החימוש בנ''מ ונ''ט
		\item הכשרת כוחות צבא שיכולים להתמודד עם צה''ל בשדה הקרב
	\end{itemize}
	
	\subsection*{הקונספציה}
	
	למה ישראל הופתעה? 
	
	\begin{itemize}
		\item בגלל אופוריה ממלחמת ששת הימים, ששררה בצבא ובדרג המדיני, והקרינה לעם. קונספציה=הנחה. אומנם, אי־אפשר לעשות מודיעין בלי הנחות יסוד, אי אפשר לעשות ישיבות על כל שידור והאזנה. העובדה שיש הנחות בסיס, זה בסדר. הבעיה, היא שכאשר הקונספציה/ההנחות חזקות מדי, קשה לקרוא נכון את המפה. בהקשר של מלחמת יום הכיפורים, גם כאן היה מודיעין. הידיעות לא פורשו נכון. היה קצין מודיעין שהודיע שבתצלומי האוויר רואים כוחות מצריים זזים. זה פשוט לא הגיע לשום מקום. ההנחות היו משהו כזה: \textit{לא תפרוץ מלחמה עד 1975}. למה? כי מצריים לא תהיה מוכנה, כי היא צריכה נ''ט ונ''מ, וסוריה לא מסוגלת ולא תצא למלחמה בלי מצריים. לכן, כאשר מזיזים כוחות באזור הדרום, תחת ההנחה הזו, הזזת כוחות יכולה להתפרש כתרגיל ולא כמלחמה. גם מצרים וסוריה עשו תרגילי הטעייה, מודיעות כמה חודשים לפני שהולך להיות תרגיל, וכו'. ''הם עשו לנו תרגיל``. 
		
		\item למדינת ישראל אין עומק אסטרטגי, היא מדינה צרה וארוכה. בגלל זה נבנה מערך מודיעין משוכלל שאמור לתת 48 שע' התראה לפני מלחמה (זמן התארגנות מילואים). זו נקודת מוצא מבינת הצבא. היו בין היתר אמצעיים מיוחדים בגבול מצריים. ישראל המשיכה לבנות חיל אויר חזק יותר לאחר ההצלחה בששת הימים, מתוך אמונה כי חיל האוויר יוכל להכריע את המלחמה העתידית בכח העליונות האווירית. מאז, בכל המלחמות גילינו שאי אפשר לנצח מלחמה עם חיל האוויר בלבד. 
		
		\item היו אמצעי האזנה בעומק השטח המצרי. היו צריכים להחליף שם סוללות שלא הוחלפו. אז היו אמצעים, שיכלו לתת התראה, אך הם לא היו פעילים. 
		
		\item לאורך תעלת סואץ הוקם קו ביצורים, והקונספציה הייתה שקו בר־לב (קו הביצורים) יעצור אותם. רוב המוצבים בקו בר־לב היו ללא אנשים, ואלו שהיו היו עם מצבת חג. גם בשבעה באוקטובר אנשים היו במצבת חג, והיה קו ביצורים שנועד למנוע כניסה. 
		
		\item המצרים והסורים הצליחו לשמור על החלטתם לצאת למלחמה בסוד מוחלט והקפידו על מידור מלא. נשתל מידע ''מרגיע`` שאותת על הכוונה לפזר את הכוחות בתום ה''תרגיל``. המידור היה טוב. החיילים בשטח ידעו רק 6 שעות לפני פרוץ המלחמה, שעומדת לפרוץ כזו. 
	\end{itemize}
	
	כבר בראשית 1973 התגשה במודיעין ההבנה כי המלך הצפוי הבא יהיה מהלך צבאי מוגבל במטרה להוביל למהלך מדיני. ההבנה הזו לא הובילה לשינוי הקונספציה. כחודש לפני המלחמה החלו לזרום ידיעות על ההכנות המצריות והסוריות להתקפה. בגלל הקונספציה, חלקן לא הגיעו אפילו לראש אמ''ן, ואלו שהגיעו נפסלו בתואנה שהסבירות למלחמה נמוכה ביותר. כבר ב־6 לאוקטובר בבוקר יום הכיפורים, הגיעה ידיעת מודיעין אמינה על פתיחת התקפה, והממשלה החליטה לכייס את מהלך המילואים. הגיוס בוצע בחפזון והכוחות שגייסו היו מצומצמים (לוקח 48 שעות לגייס מילואים). 
	
	עד היום מתקיימים דיונים האם המרגל שנתן את ידיעת הזהב היה אמין או שנסה להטעות את מדינת ישראל. הוא נתן ידיעות רבות בעבר, אך לפני המלחמה ''היה משוכנע`` שלא תפרוץ כזו. גולדה ואמ''ן האמינו לאותו מרגל. 
	
	בימים הראשונים, אנחנו נמצאים בבלימה. גם סוריה וגם מצרים פולשות לשטח מדינת ישראל בצורה משמעותי. בגבולות יש כוחות סדיר מצומצמים. לוקח זמן עד שמגיעים כוחות המילואים והכוחות הסדירים ששהו בבית. המחיר הגדול, בהרוגים, פצועים ובשטח, שולם באותם הימים הראשונים. 
	
	כמה ימים לפני המלחמה, חוסיין מלך ירדן הזהיר את גולדה שישראל עומדת לפתוח במלחמה. המודיעין התייחס בזלזול לידיעה הזו. 
	
	באותה העת שר הבטחון היה משה דיין, הרמתכ''ל דוד אלעזר (דדו), ראש אמ''ן אלי זעירא, ואלוף פיקוד דרום גורודיש. ב־5 בחודש הרמתכ''ל החליט לתגבר בכוחות שריון. הדרג המדיני לא אישר לתת מכת מנע, לתקוף מהאוויר ומהקרקע לפני פתיחת המלחמה. בניגוד למידע שהתקבל, המלחמה נפתחה בצהריים ולא בערב. 
	
	יחסי הכוחות היו מאות ישראלים לעומת עשרות אלפי ואף מאות אלפי חיילים ממדינות ערב. 
	
	הערכת המודיעין הייתה כי השריון הסורי יתקוף מצפון הרמה, שם צה''ל חיזק כוחות ואף הצליח לחסום את התקדמות הסורים. אך הסורים תקפו מצפון הרמה והתקדמו לעבר הכנרת. חלק מכוחות המילואים הגיעו בלילה הראשון. ביו םהשני ביחד עם חיל האוויר הסורים נבלמו. ביום השלישי למלחמה הפיקוד עבר להתקפת נגד. עד היו החמישי ללחימה נהדפו הסורים ורק החרמון נותר. 
	
	ביום השישי הפיקוד עבר לתקוף בתוך שטחי סוריה ותקפו מובלעת בתוך שטחי סוריה, והגיעו לטווח תותחים מדמשק. כוח ירדני ועירקי שהתווספו שינו את יחסי הכוחות וישראל החלה להגן על המובלעת שנכבשה. 
	
	האבדות הכבדות שנגרמו למצרים הובילו להחלטה לתקוף. בחרו לתקוף באיזשהו אגם אקראי באמצע הכוחות המצריים בתעלת סואץ. באיזור שנקרא החווה הסינית, כוחות צה''ל נתקלו בכוחות מצריים גדולים. באותה העת חטיבה אחרת צלחה את התעלה בסירות גומי. ביום שלאחר מכן צנחנים נוספים ניסו לתקוף בחווה הסינית ונכשלו. בשלב הזה כבר היו גשרים על התעלה וב־19 בחודש הכוחות הצליחו להגיע לכביש סואץ־קהיר. 
	
	ב־22 בחודש צה''ל כבש מידי הסורים את מוצבי החרמון שנכבשו בתחילת המלחמה. בחזית הצפון בגלל הקרבה לדמשק נכנס הלתקופה הפסקת אש. באופן דומה במצריים ב־24 באוקטובר, לאחר כיתור הארמיה השלישית. 
	
	אין צורך לדעת את מהלכי המלחמה המדויקים. זה לא חשוב לבגרות. גולדה התעקשה על דברים רבים עד להחזרת השבויים (כדוגמת אספקת אוכל ומים לארמיה השלישית). 
	
	איך השפיעה הקונספציה על הימים הראשונים? 
	\begin{multicols}{2}
		\begin{itemize}
			\item יתרון הפתעה למדינות ערב
			\item לא היו מחסני חירום מוכנים, שרר בלבול בדרג הצבאי והמדיני, בשטח לא היו כוחות
			\item הגיוס היה מאוחר וחלקי
			\item בימים הראשונים נערכו קרבות בלימה קשים
		\end{itemize}
	\end{multicols}
	
	המלחמה כולה ערכה 18 ימים. 
	
	
	\subsection*{תוצאות}
	תוצאות, נחשבות כל מה שהיה בטווח המידי, מיד לאחר המלחמה. כל דבר מעבר לשניה אחר המלחמה, נחשב השפעות. 
	
	התוצאות הברורות: 2656 הרוגים, 7251 פצועים, 301 שבויים, עשרות נעדרים (לא את כולם אתרו), אבדות גדולות בציוד צבאי במטוסים ובטנקים, ומחיר כלכלי גדול. למרות המחיר הגדול, נצחנו את המלחמה. חדרו 41km מדמשק ו־101km מקהיר. 
	
	מחר נדבר על ההשפעות. 
	
	הלוויות נערכו אחרי המלחמה ברובן מחמת כמות ההרוגים הגדולה. 
	
\end{document}